\begin{abstract}
Urban transit operators need stop-level short-term demand forecasts to
dispatch vehicles, adjust schedules, and allocate drivers in real time.
We study 3-class (low/medium/high) ridership demand classification on
the ZTBus dataset (1,409 trolleybus missions, Zurich 2019--2022) under
mission-level temporal splits (pre-2022 train, 2022 test) that mirror
realistic deployment. Our best model, a Random Forest with stop-level
lag features, reaches temporal macro-F1 = 0.8287 ($+0.2221$ over a
57-feature cross-sectional baseline of GPS, speed, power, and calendar
features). Along the way we identify and fix a subtle target-leakage
pattern that arises when standard 1~Hz lag features are computed on
forward-filled sparse passenger counts: the naive lag becomes a
near-exact copy of the current value for roughly 95\% of rows.
Computing lag features at the natural stop-event granularity, then
forward-filling the lag values, eliminates the leakage while
preserving the autoregressive signal. We additionally report a
negative result: adding TimesFM~2.5-200M penultimate-layer embeddings
(1280-dim, PCA-reduced to 32) on top of stop-level lags does not
improve any model (F1 = 0.8241, a change of $-0.005$ relative to
stop-lags-only under a single fixed training seed). The 32 retained
components capture 100\% of the training-embedding variance,
suggesting the foundation model's bus-ridership representations
collapse onto a low-rank subspace already spanned by recent-stop
history. Our contributions are (i) a strong applied result for
stop-level transit demand classification, (ii) a diagnosis and
practical fix for the forward-fill lag trap as a general
practitioner's caution, and (iii) a null result on zero-shot
foundation-model feature extraction for structured event data.
\end{abstract}
